Le gouvernement, les administrations publiques, les entreprises, les médias, les élus, les associations, et plus largement les populations, sont désormais davantage sensibilisés aux crises et à la notion de « résilience », largement diffusée dans le débat public à la suite de la crise sanitaire liée à la COVID-19\footcite{Paries2020_resilience_covid19}, et qui désigne désormais, dans l'usage courant, « la capacité à bien encaisser les coups, à surmonter les traumatismes\footcite[1]{Paries2020_resilience_covid19} ».

Une crise est qualifiée de majeure lorsque « l'étendue des phénomènes qui la caractérisent et l'intensité des troubles et transformations qui en résultent conduisent à des pertes et des dommages socialement inacceptables\footcite{info_pompiers_crise_majeure_def} ». Les seuils associés à l'étendue, à l'intensité et au degré d'acceptabilité sociale d'une crise relèvent d'une appréciation à la fois technique (quantité de victimes, perte matérielle, etc.) et politique (réaction des populations, implication et positionnement des pouvoirs publics)\footcite{info_pompiers_crise_majeure_def}.

Les crises majeures auxquelles les services de l'État sont désormais confrontés prennent des formes plus diverses et plus imbriquées : sanitaires, environnementales, technologiques, sécuritaires\footcite{AN:2022:Resilience_nationale} ou encore informationnelles\footcite[3]{SciencesPo:2025:Resilience_menaces_informationnelles}.
Cette évolution se manifeste notamment par le développement des formes de conflictualité dites « hybrides\footcite{tenenbaum2016guerre} » et dans l'amplification des phénomènes de désinformation par les réseaux sociaux\footcite{garrigou2025desinformation}. Les frontières traditionnelles entre temps de paix, temps de crise et confrontation ouverte s'en trouvent moins nettes, ce qui rend plus difficile la seule réponse par des dispositifs institutionnels spécialisés.

La planification et la mise en place de solutions de gestion de crise ne peuvent donc plus être envisagées exclusivement par le prisme des autorités compétentes. Elles supposent également une meilleure articulation avec les ressources présentes au sein de la société civile, en particulier lorsqu'elles peuvent être identifiées, qualifiées et sollicitées dans un cadre clair\footcite{AN:2022:Resilience_nationale}.

\clearpage

Cette préoccupation rejoint le point de vue de la Commission européenne qui, le 19 mars 2025\footcite{burgdorff2025livreblanc}, a dévoilé, dans un contexte marqué par la menace russe, un livre blanc pour la défense européenne destiné à définir le périmètre du plan ReArm Europe à l'horizon 2030\footcite{viepublique2025livreblanc}. Au-delà de sa dimension capacitaire, cette dynamique rappelle que la préparation aux crises engage aussi la cohésion, l'information et la capacité des sociétés à comprendre les situations auxquelles elles sont exposées.

Dans ce prolongement, les initiatives engagées au niveau national traduisent également une volonté de mieux associer la société à la préparation des crises. Le projet de service national, tel que présenté par le gouvernement français en 2026, s'inscrit dans cette dynamique en visant à impliquer davantage les jeunes générations dans des missions utiles à la défense nationale.

Les appelés du service national sont ainsi amenés à développer des savoir-faire reconnus au travers de missions relevant de trois grands domaines : la protection du territoire, le soutien au fonctionnement quotidien des armées et l'apport d'expertises dans des domaines techniques ou spécialisés\footcite{info_gouv_service_national}.

Toutefois, ce dispositif ne s'inscrit pas dans une logique de retour à une conscription généralisée, comme l'a rappelé le général Pierre Schill, chef d'état-major de l'armée de Terre\footcite{leparisien_schill_service_national}. Il traduit plutôt une évolution des modalités d'engagement, fondée sur une participation ciblée et adaptée aux besoins contemporains.

C'est dans ce cadre que les échanges relatifs aux travaux du comité se sont appuyés sur l'hypothèse de travail suivante : les \AssociationsUnionIHEDN, et en particulier l'\AssociationRegionaleIDF{} (\ARSeize), peuvent contribuer utilement à cette réflexion, à condition de préciser la nature réelle de cet apport, ses limites et ses modalités possibles. Le présent rapport de comité s'attache ainsi à examiner la problématique suivante :

\begin{quote}
\textbf{Crises majeures : quel rôle pour les citoyens engagés et les associations IHEDN ?}
\end{quote}

\clearpage

L'objectif de ce travail est d'analyser la capacité de contribution des auditeurs IHEDN face aux crises majeures, d'identifier leur plus-value spécifique et de distinguer les formes d'engagement les plus pertinentes, notamment en amont de la crise et dans les démarches de retour d'expérience. Cette analyse conduit également à poser la question d'une connaissance plus précise des profils, des compétences et des disponibilités des membres, condition nécessaire à toute contribution lisible et réaliste. L'exploitation de l'annuaire 2022 apporte ici un premier cadrage : elle permet d'estimer le périmètre associatif disponible, mais confirme aussi qu'un annuaire ne suffit pas à apprécier l'engagement effectif ni les capacités mobilisables.

Dans la suite du rapport, l'expression « \AssociationsUnionIHEDN{} » désigne l'ensemble des associations affiliées à l'\UnionIHEDN{} ; l'expression « \AssociationsRegionalesIHEDN{} » vise plus spécifiquement les associations régionales, dont l'\AssociationRegionaleIDF{} (\ARSeize). Le terme « communauté IHEDN » renvoie, pour sa part, à l'ensemble plus large des anciens auditeurs et membres associés recensés, qu'ils soient ou non adhérents d'une \AssociationUnionIHEDN.